\documentclass{blog}

\title{Analysis}

\begin{document}
\maketitle
\tableofcontents

\section{Introduction}
\begin{definition}[analysis]
    Analysis of real numbers, sequences and series of real numbers, and real-valued functions
    \cite{TAO_2022_ANALYSIS_1, TAO_2022_ANALYSIS_2}.
\end{definition}

\section{Natural Number}
\begin{definition}[natural number]
    A natural number is any element of the set
    \begin{equation}
        \mathbb{N} := {0, 1, 2, 3, \dots}
    \end{equation}
\end{definition}

\begin{axiom}[0]
    \label{axiom:peano-axiom-1}
    0 is a natural number.
\end{axiom}

\begin{axiom}[$n{++}$]
    \label{axiom:peano-axiom-2}
    If $n$ is a natural number, then $n{++}$ is also a natural number.
\end{axiom}

\begin{definition}
    \label{def:nat}
    $1 := 0{++}$, $2 := 1{++}$, $3 := 2{++}$, $\dots$
\end{definition}

\begin{axiom}[successor]
    \label{axiom:peano-axiom-3}
    0 is not the successor of any natural number,
    i.e. $n{++} \neq 0$ for every natural number $n$.
\end{axiom}

\begin{lemma}
    $4 \neq 0$
\end{lemma}

\begin{proof}
    By definition \ref{def:nat}, $3{++} = 4$.

    By Proposition \ref{prop:3-is-nat}, 3 is a natural number.

    By Axiom \ref{axiom:peano-axiom-3}, $3{++} = 4 \neq 0$.
\end{proof}

\begin{axiom}
    \label{axiom:peano-axiom-4}
    $n \neq m \Rightarrow n{++} \neq m{++}$ or $n{++} = m{++} \Rightarrow n = m$
\end{axiom}

\begin{axiom}[principle of mathematical induction]
    \label{axiom:peano-axiom-5}
    Suppose that $P(0)$ is true, and suppose that whenever $P(n)$ is true, $P(n{++})$ is also true. Then $P(n)$ is true for every natural number $n$.
\end{axiom}

\subsection{Addition}
\begin{definition}[Addition]
    Let $n$ and $m$ be natural numbers.
    \begin{enumerate}
        \item $0+m:=m$
        \item $(n{++})+m:=(n+m){++}$
    \end{enumerate}
\end{definition}

\begin{lemma}
    Let $n$ and $m$ be natural numbers. Then $n+m$ is also a natural number.
\end{lemma}

\begin{lemma}
    Let $n$ and $m$ be natural numbers.
    \begin{enumerate}
        \item $n+0=n$
        \item $n+(m{++})=(n+m){++}$
    \end{enumerate}
\end{lemma}

\begin{lemma}
    Let $a$, $b$, $c$ be natural numbers.
    \begin{enumerate}
        \item $a+b=b+a$
        \item $(a+b)+c=a+(b+c)$
        \item $a+b=a+c \Rightarrow b=c$
    \end{enumerate}
\end{lemma}

\begin{definition}
    A natural number $n$ is positive iff (if and only if) $n \neq 0$.
\end{definition}

\begin{lemma}
    \label{prop:positive-natural-numbers}
    Let $a$ and $b$ be natural numbers.
    \begin{enumerate}
        \item \label{item:positive-natural-numbers-1} If $a$ is positive, then $a+b$ is positive.
        \item \label{item:positive-natural-numbers-2} $a+b=0 \Rightarrow a=b=0$.
        \item \label{item:positive-natural-numbers-3} If $a$ is positive, then there exists only one natural number $b$ such that $b{++}=a$.
    \end{enumerate}
\end{lemma}

\begin{proof}
    Prop. \ref{prop:positive-natural-numbers} (\ref{item:positive-natural-numbers-1}).

    $a+0=a$ is positive.
    Suppose that $a+b$ is positive.
    Then $a+b{++}=(a+b){++} \neq 0$ is positive.
\end{proof}

\begin{proof}
    Prop. \ref{prop:positive-natural-numbers} (\ref{item:positive-natural-numbers-3}).

    Suppose that $a=b{++}$.
    Then $a{++}=(a){++}$.
\end{proof}

\begin{definition}
    Let $n$ and $m$ be natural numbers.
    \begin{enumerate}
        \item $n \geq m$ or $m \leq n$ iff $n=m+a$
        \item $n>m$ or $m<n$ iff $n \geq m$ and $n \neq m$
    \end{enumerate}
\end{definition}

\begin{lemma}
    \label{prop:natural-numbers-order}
    Let $a$, $b$, $c$ be natural numbers.
    \begin{enumerate}
        \item \label{item:natural-numbers-order-1} $a \geq a$
        \item \label{item:natural-numbers-order-2} If $a \geq b$ and $b \geq c$, then $a \geq c$
        \item \label{item:natural-numbers-order-3} If $a \geq b$ and $b \geq a$, then $a=b$
        \item \label{item:natural-numbers-order-4} $a \geq b$ iff $a+c \geq b+c$
        \item \label{item:natural-numbers-order-5} $a<b$ iff $a{++} \leq b$
        \item \label{item:natural-numbers-order-6} $a<b$ iff $b=a+d$ for some positive number $d$
    \end{enumerate}
\end{lemma}

\begin{proof}
    Prop. \ref{prop:natural-numbers-order} (\ref{item:natural-numbers-order-5}).

    $b=a+d$ and $b \neq a \Rightarrow d \neq 0 \Rightarrow d = n{++}$

    $b=a+(n{++})=(a{++})+n \Rightarrow a{++} \leq b$
\end{proof}

\begin{lemma}
    Let $a$ and $b$ be natural numbers. Then exactly one of following statement is true:
    $a<b$, $a=b$, or $a>b$.
\end{lemma}

\begin{proof}
    If $a<b$ and $a>b \overset{\text{Prop. \ref{prop:natural-numbers-order} (\ref{item:natural-numbers-order-3})}}{\Rightarrow} a=b$ contradiction
\end{proof}

\begin{lemma}[Strong principle of induction]
    Let $m_{0}$ be a natural number. Suppose that for each $m \geq m_{0}$, if $P(m')$ is true for all $m_{0} \leq m' < m$, then $P(m)$ is also true. Then $P(m)$ is true for all $m \geq m_{0}$.
\end{lemma}

\begin{proof}
    Let $Q(n)$ be the property that $P(m)$ is true for all $m_{0} \leq m < n$.

    Suppose that if $Q(n)$ is true, then $P(m)$ is also true.

    Then $P(m)$ is true for all $m_{0} \leq m < n{++}$.
\end{proof}

\begin{example}
    Fibonacci numbers.
    \begin{gather}
        \label{eq:fibonacci-numbers}
        F_{n} = \frac{\phi^{n} - \psi^{n}}{\phi-\psi} \\
        F_{n+2} = F_{n+1} + F_{n}
    \end{gather}
    where $\phi=\frac{1+\sqrt{5}}{2}$ and $\psi=\frac{1-\sqrt{5}}{2}$
\end{example}

\begin{example}
    Principle of backwards induction.

    Let $n$ be a natural number. Whenever $P(m{++})$ is true, $P(m)$ is true.
    If $P(n)$ is true, then $P(m)$ is true for all $m \leq n$.
\end{example}

\begin{proof}
    Suppose that if $P(n)$ is true, then $P(m)$ is true for all $m \leq n$.
    If $P(n{++})$ is true, then $P(m)$ is true for all $m \leq n{++}$.
\end{proof}

\begin{example}
    Principle of induction starting from the base case $n$.

    Let $n$ be a natural number. Whenever $P(m)$ is true, $P(m{++})$ is true.
    If $P(n)$ is true, then $P(m)$ is true for all $m \geq n$.
\end{example}

\begin{proof}
    If $P(0)$ is true, suppose that $P(m)$ is true.
    Then $P(m{++})$ is also true.
    Then $P(m)$ is true for all $m \geq 0$.
    Suppose that if $P(n)$ is true, then $P(m)$ is true for all $m \geq n$.
    If $P(n{++})$ is true, then $P(m)$ is true for all $m \geq n{++}$.
\end{proof}

\subsection{Multiplication}
\begin{definition}
    Let $n$ and $m$ be natural numbers.
    \begin{enumerate}
        \item $0 \times m := 0$
        \item $(n{++}) \times m := (n \times m) + m$
    \end{enumerate}
\end{definition}

\begin{lemma}
    Let $n$ and $m$ be natural numbers. Then $n \times m$ is also a natural number.
\end{lemma}

\begin{lemma}
    Let $n$ and $m$ be natural numbers.
    \begin{enumerate}
        \item $n \times 0 = 0 \times n$
        \item $n \times (m{++}) = (n \times m) + n$
    \end{enumerate}
\end{lemma}

\begin{lemma}[Commutative]
    \label{lemma:multiplication-commutative}
    Let $n$ and $m$ be natural numbers.
    \begin{enumerate}
        \item $n \times m = m \times n$
        \item $n \times m = 0$ iff at least one of $n$ and $m$ is equal to 0.
              If $n$ and $m$ are both positive, then $nm$ is also positive.
    \end{enumerate}
\end{lemma}

\begin{lemma}
    Let $a$, $b$, $c$ be natural numbers.
    \begin{enumerate}
        \item $a(b+c)=(b+c)a=bc+ca$
        \item $(ab)c=a(bc)$
        \item If $a<b$ and $c$ is positive, then $ac<nc$.
        \item If $ac=bc$ and $c$ is positive, then $a=b$.
    \end{enumerate}
\end{lemma}

\begin{lemma}[Euclid's division lemma]
    Let $n$ be a natural number and $q$ be a positive natural number. Then there exist natural numbers $m$ and $r$, such that $0 \leq r < q$ and $n=mq+r$.
\end{lemma}

\begin{proof}
    $0 = 0 \times q + 0$ and $0 \leq r < q$.
    Suppose that $n=mq+r$ and $0 \leq r < q$.
    Then $n{++}=mq+(r{++})$ and $0 \leq r{++} \leq q$.
    If $r{++}=q$, $n{++}=mq+q=(m+1)q+0$.
    If $r{++}<q$, $n{++}=mq+r{++}$.
\end{proof}

\begin{definition}[Exponentiation]~
    \begin{enumerate}
        \item $m^{0}:=1$ and $0^{0}:=1$
        \item $m^{n{++}}:=m^{n} \times m$
    \end{enumerate}
\end{definition}

\section{Set Theory}
\begin{definition}
    Set $A$ is any unordered collection of objects.
    $x \in A$ if $x$ lies in the collection, otherwise $x \notin A$.
\end{definition}

\begin{axiom}~
    \label{axiom:sets}
    \begin{enumerate}
        \item \label{item:sets-1} If $A$ is a set, then $A$ is also an object.
        \item \label{item:sets-2} $A=B$ iff every element of $A$ belongs also to $B$ and vice versa.
        \item \label{item:sets-3} There exists an empty set $\emptyset$, such that every object $x \notin \emptyset$.
        \item \label{item:sets-4} Let $A$ be a non-empty set. Then there exist an object $x \in A$.
        \item \label{item:sets-5} Singleton set $\{a\}$: $y \in \{a\}$ iff $y=a$.
              Pair set $\{a,b\}$: $y \in \{a,b\}$ iff $y=a$ or $y=b$.
        \item \label{item:sets-6} $x \in A \cup B \Leftrightarrow x \in A$ or $x \in B$.
        \item \label{item:sets-7} $y \in \{x \in A : P(x) \text{ is true}\} \Leftrightarrow y \in A \text{ and } P(x) \text{ is true}$.
        \item \label{item:sets-8} $z \in \{y : P(x,y) \text{ is true for some } x \in A\} \Leftrightarrow P(x,z) \text{ is true for some } x \in A$.
        \item \label{item:sets-9} There exists a set $\mathbb{N}$, whose elements are called natural numbers, and an object $n{++}$ assigned to every $n \in \mathbb{N}$, such that the Peano axioms \ref{axiom:peano-axiom} hold.
    \end{enumerate}
\end{axiom}

\begin{example}
    $\{a,a\} = \{a\}$
\end{example}

\begin{example}
    $\emptyset$, $\{\emptyset\}$, $\{\{\emptyset\}\}$, $\{\emptyset,\{\emptyset\}\}$ are all distinct.
\end{example}

\begin{proof}
    $y \notin \emptyset$.

    $y=\emptyset \in \{\emptyset\}$.

    $y=\{\emptyset\} \in \{\{\emptyset\}\}$.

    $y=\emptyset$ or $y=\{\emptyset\} \in \{\emptyset,\{\emptyset\}\}$.
\end{proof}

\begin{lemma}
    Let $a$, $b$ are objects, and $A$, $B$, $C$ are sets.
    \begin{enumerate}
        \item $\{a,b\}=\{a\}\cup\{b\}$
        \item $A \cup B = B \cup A$
        \item $(A \cup B) \cup C = A \cup (B \cup C)$
    \end{enumerate}
\end{lemma}

\begin{definition}
    \begin{gather*}
        \{a,b,c\} := \{a\}\cup\{b\}\cup\{c\} \\
        \{a,b,c,d\} := \{a\}\cup\{b\}\cup\{c\}\cup\{d\} \\
        \vdots
    \end{gather*}
\end{definition}

\begin{definition}[Subsets]
    \begin{enumerate}
        \item $A \subseteq B \Leftrightarrow$ for any object $x \in A$, then $x \in B$.
        \item $A \subsetneq B$ if $A \subseteq B$ and $A \neq B$.
    \end{enumerate}
\end{definition}

\begin{lemma}
    Let $A$, $B$, $C$ be sets.
    \begin{enumerate}
        \item If $A \subseteq B$ and $B \subseteq C$, then $A \subseteq C$.
        \item If $A \subseteq B$ and $B \subseteq A$, then $A = C$.
        \item If $A \subsetneq B$ and $B \subsetneq C$, then $A \subsetneq C$.
    \end{enumerate}
\end{lemma}

\begin{definition}
    Let $A$, $B$ be sets.
    \begin{enumerate}
        \item $A \cap B := \{x \in A : x \in B\}$.
        \item $A \backslash B := \{x \in A : x \notin B\}$.
    \end{enumerate}
\end{definition}

\begin{lemma}
    \label{prop:sets-boolean-algebra}
    Let $A$, $B$, $C$ be sets, and $\{A,B,C\} \subseteq X$.
    \begin{enumerate}
        \item \label{item:sets-boolean-algebra-1} $A \cup \emptyset = A$ and $A \cap \emptyset = \emptyset$.
        \item \label{item:sets-boolean-algebra-2} $A \cup X = X$ and $A \cap X = A$.
        \item \label{item:sets-boolean-algebra-3} $A \cup A = A$ and $A \cap A = A$.
        \item \label{item:sets-boolean-algebra-4} $A \cup B = B \cup A$ and $A \cap B = B \cap A$.
        \item \label{item:sets-boolean-algebra-5} $(A \cup B) \cup C = A \cup (B \cup C)$ and $(A \cap B) \cap C = A \cap (B \cap C)$.
        \item \label{item:sets-boolean-algebra-6} $A \cap (B \cup C) = (A \cap B) \cup (A \cap C)$ and $A \cup (B \cap C) = (A \cup B) \cap (A \cup C)$.
        \item \label{item:sets-boolean-algebra-7} $A \cup (X \backslash A) = X$ and $A \cap (X \backslash A) = \emptyset$.
        \item \label{item:sets-boolean-algebra-8} De Morgan laws: $X \backslash (A \cup B) = (x \backslash A) \cap (X \backslash B)$ and $X \backslash (A \cap B) = (X \backslash A) \cup (X \backslash B)$.
    \end{enumerate}
\end{lemma}

\begin{proof}
    Prop. \ref{prop:sets-boolean-algebra} (\ref{item:sets-boolean-algebra-5}).

    $x \in A \cap X \Leftrightarrow x \in A \Rightarrow A \cap X = A$.
\end{proof}

\begin{proof}
    Prop. \ref{prop:sets-boolean-algebra} (\ref{item:sets-boolean-algebra-6}).

    $x \in A \cap (B \cup C) \Rightarrow x \in A \text{ and } (x \in B \text{ or } x \in C) \Rightarrow x \in A \cap B \text{ or } x \in A \cap C \Rightarrow x \in (A \cap B) \cup (A \cap C)$.

    $x \in (A \cap B) \cup (A \cap C) \Rightarrow x \in (A \cap B) \text{ or } x \in (A \cap C) \Rightarrow x \in A \text{ and } x \in B \cup C \Rightarrow x \in A \cap (B \cup C)$.
\end{proof}

\begin{lemma}[De mongan laws]~

    !(A or B) = !A and !B.

    !(A and B) = !A or !B.
\end{lemma}

\begin{proof}
    Prop. \ref{prop:sets-boolean-algebra} (\ref{item:sets-boolean-algebra-8}).

    $x \in X \backslash (A \cup B) \Rightarrow x \in X \text{ and } (x \notin A \text{ and } x \notin B)$.

    $x \in X \backslash (A \cap B) \Rightarrow x \in X \text{ and } (x \notin A \text{ or } x \notin B)$.
\end{proof}

\begin{example}~

    $\{y : y=f(x) \text{ for some } x \in A\} = \{f(x) : x \in A\}$.

    $\{x \in A : P(x) \text{ is true}\} \Rightarrow \{f(x) : x \in A; P(x) \text{ is true}\}$.
\end{example}

\begin{example}
    $\{n+3 : n \in \mathbb{N}, 0 \leq n \leq 5\} = \{8-m : m \in \mathbb{N}, 0 \leq m \leq 3\}$.
\end{example}

\begin{example}[Axiom \ref{axiom:sets} (\ref{item:sets-8}) $\Rightarrow$ Axiom \ref{axiom:sets} (\ref{item:sets-7})]~

    $y \in \{x:P(x) \text{ is true for some } x \in A\} \Rightarrow y \in \{x \in A : P(x) \text{ is true}\}$.
\end{example}

\subsection{Functions}

\section{Power Series}

\printbibliography
\end{document}
